% =====================================================================
%  AI4Good workshop variant.
%  Defines everything main.tex pulls in: title, venue string, abstract,
%  introduction, discussion. Nothing else in the paper changes.
% =====================================================================

\newcommand{\wkTitle}{Portable Causal Fairness Across \\
Synthetic Data Generator Families}

\newcommand{\wkVenue}{AI for Good}

% ---------------------------------------------------------------- abstract
\newcommand{\wkAbstract}{%
When a statistical agency or regulator releases synthetic data in place of
sensitive records, it chooses the generator that produces the table, and can
shape that generator so unfair pathways are absent. DECAF made this concrete
on one non-private GAN: three fairness definitions become three sets of edge
cuts on the generator's causal graph. Whether the mechanism belongs to DECAF,
or to causal factorisation itself, was untested. We port all three definitions
to nine generators from three unrelated families (marginals-based, GAN, and
diffusion, each with differentially private variants), across three levels of
formal privacy guarantee, over 2{,}520 matched-pair runs on Adult and COMPAS.
The mechanism transfers everywhere, and our new causal diffusion backbone
yields the fairest release of any family we tested, at fidelity close to the
marginals tier. Applying the cut barely moves fidelity, and its cost in
downstream classifier accuracy is small (about $0.07$ to $0.15$ AUC on average,
depending on how aggressive the cut is).%
}

% ------------------------------------------------------------------- intro
\newcommand{\wkIntro}{%
Institutions holding sensitive records increasingly publish a synthetic
substitute rather than the records themselves: a hospital, a lending regulator
or a statistical agency fits a generative model and releases a table drawn from
it. The practice is real and growing: the 2020 US Census was released under a
differential-privacy guarantee \citep{abowd2022topdown}, and public challenges
run by NIST, the FDA and the US--UK privacy-technology programme have made
private tabular release a working technology rather than a proposal
\citep{nist2018challenge,V-CHAMPS,vos2024}. That substitution is usually justified on privacy grounds, but it carries a
second, less-exploited consequence. Whoever generates the data chooses the
generating process, and a process can be edited. A disparity that exists in the
world because of a pathway running from a protected attribute, such as national
origin, to an outcome, such as whether a loan is approved, need not exist in
the released table.

The cleanest existing realisation of that lever is
DECAF~\citep{vanbreugel2021decaf}, which factors a generator along a causal
graph so that severing an edge at sampling time makes the generator
structurally unable to use the pathway it named. But DECAF is a single
non-private WGAN-GP~\citep{gulrajani2017wgangp}, whose released tables trail
modern differentially private marginals-based synthesizers and modern
diffusion models on the fidelity and utility benchmarks that decide what
institutions actually deploy (Sec.~\ref{sec:results}). Whether the fairness mechanism is confined to
that one architecture, or is a property of causal factorisation itself, is
therefore a practical question and not merely an academic one: it decides whether an
institution deploying synthetic data has the intervention available at all.

We ask two questions of that mechanism. Is edge surgery a property of DECAF,
or of the causal factorisation any such generator uses? And can it be attached
to a generator whose graph is not a supplied causal DAG but the undirected
marginal-dependency graph a differentially private synthesizer builds from
noisy counts? The stakes are practical: if the mechanism that would make a
released table fair is confined to a generator no institution would actually
deploy, it might as well not exist.%
}

% -------------------------------------------------------------- discussion
\newcommand{\wkDiscussion}{%

\textbf{Edge surgery belongs to causal factorisation, not to DECAF.} DECAF
conjectured that its mechanism would extend beyond its own GAN but tested no
other generator, and PreFair implements one such definition (CF) on a single
marginals synthesizer (MST). We show the definitions compile identically
across nine generators, three per family, on both datasets and at every
privacy budget; an institution can now pick a synthesizer for its fidelity,
privacy or maturity and keep the fairness lever, instead of being pushed onto
a non-private GAN or a single marginals implementation to get it.

\textbf{Causal diffusion produces the fairest release we tested.} Averaged
over both datasets, our diffusion backbone leaves the smallest downstream
parity gap under both CF ($0.099$) and DemPar ($0.081$), ahead of causal GAN
($0.128$, $0.089$) and marginals ($0.152$, $0.136$; Table~\ref{tab:fairness}).
This is the finding we want to highlight, since the diffusion backbone is new
work: DECAF gave us one GAN and PreFair one marginals-based synthesizer, and
neither offered a diffusion route, through which the mechanism turns out to
reduce disparity most.

\textbf{Which generator should a publisher use?} \emph{Marginals-based
synthesizers} remain the default when fidelity and privacy come first: they
are truest to the real table on every fidelity measure, and are
differentially private by construction. \emph{Causal diffusion} is the
choice when the fairness lever has to bite, removing the largest fraction of
the parity gap while holding fidelity close to the marginals tier. \emph{The
causal GAN family} trails on fidelity and utility, and DP-GAN fails
outright. On Adult, CF pays $6$ AUC points to remove $41\%$ of the baseline
gap on marginals, $16$ ($78\%$) on causal GAN, and $18$ ($76\%$) on causal
diffusion; the cost is smaller on COMPAS, where every family starts closer
to chance, and DemPar costs roughly double CF on both COMPAS families. A
publisher should reach for CF unless there is a specific reason to route out
every admissible pathway too.

\textbf{Only the in-generator cut ships with the table.} On datasets where
the private structure encodes little to sever, a downstream mitigator can
reach a lower parity gap on the release (Table~\ref{tab:fairness}). But a
public table is used by consumers its publisher will never meet, and of the
four places to intervene, the in-generator cut is the only one that travels
with the release and cannot be dropped.

\textbf{Conclusion.} The point is not any generator we ported to, but that
all three families share the same abstraction (a causal factorisation whose
edges can be severed), and this abstraction, not any one architecture, is
where the fairness lever lives. An institution can now pick a synthesizer for
its privacy or data-quality needs, and keep a portable causal-fairness
intervention as part of that choice.%
}
